\section{Wyniki optymalizacji}

W tym rozdziale należy przedstawić wyniki uzyskane dla wszystkich
analizowanych kombinacji architektury i współczynnika uczenia.

\subsection{Przebieg procesu uczenia}

Dla każdej konfiguracji należy przeanalizować przebieg funkcji kosztu
w kolejnych epokach.

\begin{figure}[H]
    \centering
    \fbox{
        \parbox{0.9\textwidth}{
            \centering
            Miejsce na wykresy przebiegu funkcji kosztu dla wszystkich
            analizowanych konfiguracji lub dla konfiguracji reprezentatywnych.
        }
    }
    \caption{Przebieg funkcji kosztu w procesie uczenia dla wybranych konfiguracji}
    \label{fig:optimization_loss_curves}
\end{figure}

\subsection{Porównanie metryk końcowych}

Po zakończeniu treningu dla każdej konfiguracji należy zestawić końcowe
wartości najważniejszych metryk jakości.

\begin{table}[H]
    \centering
    {\scriptsize
\setlength{\tabcolsep}{4pt}
\begin{tabular}{@{}p{0.31\textwidth}cccc@{}}
\toprule
Konfiguracja & LR & Loss & Sensitivity & Specificity \\
\midrule
\(13 \rightarrow 4 \rightarrow 3\) & \(0.001\) & \(0.7688 \pm 0.1129\) & \(0.6556 \pm 0.1127\) & \(0.8271 \pm 0.0578\) \\
\(13 \rightarrow 4 \rightarrow 3\) & \(0.01\) & \(0.0674 \pm 0.0428\) & \(0.9743 \pm 0.0164\) & \(0.9864 \pm 0.0089\) \\
\(13 \rightarrow 4 \rightarrow 3\) & \(0.1\) & \(0.0894 \pm 0.0631\) & \(0.9739 \pm 0.0186\) & \(0.9863 \pm 0.0096\) \\
\(13 \rightarrow 8 \rightarrow 3\) & \(0.001\) & \(0.5762 \pm 0.1325\) & \(0.8333 \pm 0.1724\) & \(0.9145 \pm 0.0865\) \\
\(13 \rightarrow 8 \rightarrow 3\) & \(0.01\) & \(0.0659 \pm 0.0408\) & \(0.9710 \pm 0.0209\) & \(0.9849 \pm 0.0104\) \\
\(13 \rightarrow 8 \rightarrow 3\) & \(0.1\) & \(0.0824 \pm 0.0590\) & \(0.9706 \pm 0.0268\) & \(0.9848 \pm 0.0131\) \\
\(13 \rightarrow 13 \rightarrow 3\) & \(0.001\) & \(0.4118 \pm 0.0811\) & \(0.9333 \pm 0.0971\) & \(0.9660 \pm 0.0437\) \\
\(13 \rightarrow 13 \rightarrow 3\) & \(0.01\) & \(0.0605 \pm 0.0420\) & \(0.9743 \pm 0.0164\) & \(0.9864 \pm 0.0089\) \\
\(13 \rightarrow 13 \rightarrow 3\) & \(0.1\) & \(0.0797 \pm 0.0571\) & \(0.9763 \pm 0.0204\) & \(0.9875 \pm 0.0105\) \\
\(13 \rightarrow 13 \rightarrow 8 \rightarrow 3\) & \(0.001\) & \(1.0894 \pm 0.0000\) & \(0.3333 \pm 0.0000\) & \(0.6667 \pm 0.0000\) \\
\(13 \rightarrow 13 \rightarrow 8 \rightarrow 3\) & \(0.01\) & \(0.5584 \pm 0.2835\) & \(0.7111 \pm 0.2004\) & \(0.8614 \pm 0.0948\) \\
\(13 \rightarrow 13 \rightarrow 8 \rightarrow 3\) & \(0.1\) & \(0.0648 \pm 0.0544\) & \(0.9814 \pm 0.0134\) & \(0.9905 \pm 0.0066\) \\
\(13 \rightarrow 32 \rightarrow 16 \rightarrow 3\) & \(0.001\) & \(1.0894 \pm 0.0000\) & \(0.3333 \pm 0.0000\) & \(0.6667 \pm 0.0000\) \\
\(13 \rightarrow 32 \rightarrow 16 \rightarrow 3\) & \(0.01\) & \(0.1242 \pm 0.1668\) & \(0.9443 \pm 0.1004\) & \(0.9695 \pm 0.0547\) \\
\(13 \rightarrow 32 \rightarrow 16 \rightarrow 3\) & \(0.1\) & \(0.0896 \pm 0.0721\) & \(0.9763 \pm 0.0204\) & \(0.9875 \pm 0.0105\) \\
\(13 \rightarrow 64 \rightarrow 32 \rightarrow 16 \rightarrow 3\) & \(0.001\) & \(1.0894 \pm 0.0000\) & \(0.3333 \pm 0.0000\) & \(0.6667 \pm 0.0000\) \\
\(13 \rightarrow 64 \rightarrow 32 \rightarrow 16 \rightarrow 3\) & \(0.01\) & \(1.0897 \pm 0.0000\) & \(0.3333 \pm 0.0000\) & \(0.6667 \pm 0.0000\) \\
\(13 \rightarrow 64 \rightarrow 32 \rightarrow 16 \rightarrow 3\) & \(0.1\) & \(0.1981 \pm 0.3107\) & \(0.9034 \pm 0.2027\) & \(0.9504 \pm 0.1012\) \\
\bottomrule
\end{tabular}
}

    \caption{Zbiorcze porównanie wyników dla analizowanych konfiguracji}
    \label{tab:optimization_results_summary}
\end{table}

\subsection{Czułość i specyficzność}

Zgodnie z wymaganiami projektowymi, dla każdej konfiguracji należy przedstawić
wartości czułości i specyficzności.

\begin{table}[H]
    \centering
    \begin{tabular}{lcccc}
\toprule
Konfiguracja & Klasa & Sensitivity & Specificity & Zbiór \\
\midrule
\multicolumn{5}{l}{Miejsce na uzupełnienie wartości dla analizowanych konfiguracji} \\
\bottomrule
\end{tabular}

    \caption{Czułość i specyficzność dla analizowanych konfiguracji}
    \label{tab:sensitivity_specificity_summary}
\end{table}

\subsection{Wybór najlepszej konfiguracji}


